\chapter{Conclusions and Recommendations}

\section{Conclusions}
The conclusions are stated in the same order as the six specific objectives.
\begin{enumerate}[leftmargin=1.0cm,itemsep=0.55em,topsep=0.3em]
\item \textbf{Objective 1 -- Draft Final NBC 205:2012 reference.} The representative 2012 building showed increasing concrete design-check deficiencies under the adopted NBC 105:2025 demand cases, from 1 failing frame at Soil C/$Z=0.25$ to 21 at Soil C/$Z=0.35$ and $0.40$. The result supports rechecking earlier ready-to-use buildings against the current seismic-demand framework; it is not a universal failure rate for all 2012 buildings.
\item \textbf{Objective 2 -- NBC 205:2024 Regular benchmark.} The Regular model provides the reference response for the matched comparison. At Soil C/$Z=0.40$ it recorded a Mode-1 period of 0.699~s, the largest validated EQx Story-1 shear of 532.489~kN, a governing ULS drift ratio of 0.004784, a first-peak pushover capacity of 1836.338~kN, and the largest performance-point base shear of 456.253~kN.
\item \textbf{Objective 3 -- Selected irregular configurations.} The six plan-irregular models did not show a uniform loss or amplification of response relative to the Regular model. U governed the maximum ULS drift (0.004954), first-peak pushover base shear (1882.070~kN), and maximum performance-point roof displacement (27.072~mm), whereas Non-Parallel governed the main-storey CM--CR separation (1.237~m). Other irregular models were lower than the Regular benchmark for several response measures.
\item \textbf{Objective 4 -- Plan configuration and CM--CR effects.} The complete drift profiles, detailed CM--CR coordinates, pushover summary, and performance-point comparison show that different geometric mechanisms alter different response measures. L2 and Non-Parallel have much larger CM--CR separation than the Regular reference but lower governing drift and performance-point displacement, while U remains comparatively close in EQx shear but exceeds the Regular reference in several deformation/capacity quantities. CM--CR separation is descriptive and is not used as the formal NBC torsional-irregularity ratio.
\item \textbf{Objective 5 -- Practical RUD applicability boundary.} The results do not justify unrestricted application of NBC 205:2024 RUD detailing to all irregular plans. The analysed configurations provide comparative evidence that some response checks remain within the adopted limits, but the prescriptive RUD scope and the analytical response checks are distinct. Irregular or height-extension cases therefore require project-specific engineered verification rather than automatic prescriptive acceptance. The present evidence is insufficient for a quantitative height-effect comparison of storey-force distribution because a matched four-storey response series was not verified.
\item \textbf{Objective 6 -- GII and seismic response.} GII does not produce a unique or monotonic response ranking. L and U both have GII=16.67\% yet differ in shear, drift, pushover and performance-point response; L1, L2 and T all have GII=33.33\% but also show materially different responses. GII is therefore useful as an area-deficit descriptor but is insufficient as a stand-alone seismic-performance criterion.
\end{enumerate}

\begin{itemize}[leftmargin=0.8cm,itemsep=0.35em,topsep=0.3em,label=\textbullet]
\item \textbf{Additional modal finding.} For Soil Type C at $Z=0.40$, the seven-model fundamental periods range from 0.649 to 0.707~s. Mode~1 carries approximately 79.09--80.32\% of the translational mass in the X direction, Mode~2 carries approximately 74.82--80.42\% in the Y direction, and cumulative translational participation exceeds 90\% by Mode~4 in X and Mode~5 in Y for every configuration.
\end{itemize}

\section{Recommendations}
\begin{itemize}[leftmargin=0.8cm,itemsep=0.45em,topsep=0.3em,label=\textbullet]
\item Irregular buildings should be subjected to project-specific structural analysis and design rather than treated as automatically covered by the RUD provisions for regular buildings.
\item Where torsional irregularity must be classified formally, the required diaphragm-edge displacement results should be evaluated directly in accordance with the applicable seismic code provisions.
\item A wider range of plan configurations, degrees of irregularity, building heights, soil conditions, and seismic demand levels should be investigated before proposing any extension of the prescriptive scope of NBC 205:2024.
\item Future studies should evaluate additional configurations using consistent modelling assumptions and clearly defined response measures so that the influence of geometry can be distinguished from other structural parameters.
\end{itemize}
